\documentclass[11pt]{amsart}
\usepackage{amsmath,amssymb,amsthm}
\usepackage[margin=1.05in]{geometry}
\usepackage{booktabs}
\usepackage{enumitem}
\usepackage{hyperref}
\emergencystretch=2em

\newtheorem{theorem}{Theorem}[section]
\newtheorem{lemma}[theorem]{Lemma}
\newtheorem{proposition}[theorem]{Proposition}
\newtheorem{corollary}[theorem]{Corollary}
\theoremstyle{definition}
\newtheorem{definition}[theorem]{Definition}
\newtheorem{remark}[theorem]{Remark}
\newtheorem{question}[theorem]{Question}
\newtheorem{conjecture}[theorem]{Conjecture}

\newcommand{\Ev}{\mathrm{Ev}}
\newcommand{\Od}{\mathrm{Od}}

\title{Lov\'asz's spare path: the missing move in the Fan game}
\author{Justin Williams}
\date{November 2026 (working draft)}

\begin{document}

\begin{abstract}
Fan's addible-edge transfers have hypotheses that depend only on the end profile of a path
partition. Starting from a Lov\'asz partition of $G-F$ and re-inserting $F$ by Fan transfers is
therefore a finite game on profiles. This game wins the rooted Gallai test at every root of every
connected graph of order~$6$, but it loses on some graphs of order~$8$, the smallest being a $K_4$
of even vertices with a full vertex carrying a leaf and two false twins. We show that the missing
move is not a move beyond Fan's inequality. It is \emph{path insertion}: a new path of length at
least two made of held edges, paid for by the spare path that Lov\'asz's theorem leaves when $G-F$
has an isolated vertex. The key instance is a \emph{hub lemma}. Re-inserting the four edges at an
isolated vertex costs exactly one path: two edges are received by Fan transfers, and the other
two become one new path through the hub. It yields the rooted test at every degree-four vertex all
of whose neighbours are even, provided at most one further vertex is even. A variant with one odd
neighbour closes the critical graph of order~$8$ on which the original game loses. With path
insertion
the game wins every root of every connected graph of order~$8$, with at most nine held edges, and
nine are needed for $K_8$ minus a perfect matching.
Profile-level P\'osa moves do not help. These are rotation at equality, and a strengthening of
Fan's lemma that we prove (\emph{Fan's alternative}: either the transfer succeeds, or at least
$\sum_T b$ blank neighbours of the receiver can each be given an end). Neither repairs a single
losing root. At a hub of degree six a pure star pays for one receipt too few, and the missing
receipt sits at Fan's equality. The game compensates by routing the new path along further held
edges. The rigid configuration at equality can be refuted by an edge count, or recombined as a
flower of cycles. We give a profile-level receipt-or-flower move with such a count. With it, and
the symmetry of the root's orbit, the game wins our hard cases of orders $8$ and $10$ with a
single held star. On a random sample of order $10$ every graph is won with a star plus at most
three held edges. These include the cocktail party graphs $K_{2m}$ minus a perfect matching for
$m=4,5$, where the move attains the insertion barrier. Their rooted test is already known for
every $m$, by explicit Hamiltonian flowers, so on these graphs our results are re-derivations
inside the game. For every degree we prove the hub lemma
under a supply condition, a version with relays into the hub, and a version with two hubs that
share one new path. We test a hub-strategy conjecture: at a root with the hypotheses of the
branching base case, the root's star, a parity repair and a bounded number of further held
edges suffice. At order $8$ two further edges suffice at all $9268$ root orbits. A sample of order
$10$ needs three, and there the proved theorems apply at no root. The hub lemmas of every degree,
the degree-four theorems, the critical instance and Fan's alternative are formalised in Lean~4.
\end{abstract}

\maketitle

\section{Setting and summary}\label{sec:intro}

Graphs are finite and simple. A \emph{path} is a nonempty simple path, and a \emph{partition} of a
graph is a partition of its edge set into paths. For a partition $D$ and a vertex $v$ let $b_D(v)$
be the number of path ends at $v$. Then
\begin{equation}\label{eq:parity}
b_D(v)\equiv d(v)\pmod 2,\qquad \sum_v b_D(v)=2|D|.
\end{equation}
$\Ev(G)$ and $\Od(G)$ are the sets of vertices of positive even degree and of odd degree. For a
vertex $u$ of a graph $H$ with partition $D$, the \emph{blank neighbours} of $u$ are
$Z_D(u)=\{w\in N_H(u):b_D(w)=0\}$.

\subsection*{The rooted test} Gallai's conjecture is equivalent to the following statement
(R)~\cite{ES}. \emph{For every connected graph $G$ of even order $n$ and every vertex $r$, the
graph $G$ has a partition into at most $n/2$ paths with an end at $r$.} Equivalently, $G$ plus a
pendant leaf at $r$ has at most $n/2$ paths.

\subsection*{Inputs} Only two classical facts are used.
\begin{enumerate}[label=(\Alph*)]
\item \emph{Lov\'asz~\cite{L68}.} A graph whose vertices of positive degree include at most one
of even degree has a partition with $b=d\bmod 2$. Such a partition has $|\Od|/2$ paths. If the
graph has $N$ vertices, none of them isolated, this is at most $\lfloor N/2\rfloor$. Isolated
vertices cost nothing. (Lean: \texttt{lovasz\_ends}.)
\item \emph{Fan~\cite{Fan}, Lemmas 3.3--3.4.} Let $D$ be a partition of $H$, let $u$ be a vertex,
and let $T$ be a set of non-neighbours of $u$ with $u\notin T$. If
$\sum_{t\in T}b_D(t)>|Z_D(u)|$, then for some $t\in T$ with $b_D(t)\ge1$ the graph $H+tu$ has a
partition $D'$ with $|D'|=|D|$ and $b_{D'}=b_D+e_u-e_t$. We say that $t$ \emph{gives an end to}
$u$, or that $u$ \emph{receives} from $t$. (Lean: \texttt{fan\_grouped}, \texttt{fan\_single}.)
\end{enumerate}

\subsection*{The Fan game}\label{sec:game}
Fan's hypothesis depends only on the current graph and the current profile, so the following is
a game on pairs $(H,b)$.

\begin{definition}[The Fan game]\label{def:game}
Let $G$ have even order $n$ and let $r\in V(G)$.
\begin{enumerate}
\item Choose a set $F$ of \emph{held} edges and put $K=G-F$. Start from a profile that Lov\'asz
guarantees on $K$:
\begin{itemize}
\item $b=d_K\bmod 2$ if $K$ has at most one vertex of positive even degree; or
\item $b=1$ on $\Od(K)$, $b=2$ at $v$ and $b=0$ at $u$, if $K$ has no isolated vertex and
$\Ev(K)=\{u,v\}$ (endpoint-structure Lemma~4~\cite{ES}: $\Ev(K)\setminus\{u\}$ is realisable).
\end{itemize}
\item Re-insert the held edges. Let $H$ be the current graph. The moves are the following.
\begin{itemize}
\item \emph{Single transfer} of a held edge $tu$ towards $u$, if $b(t)>|Z(u)|$.
\item \emph{Grouped transfer} at $u$ with a set $T$ of held edges $tu$, if
$\sum_Tb>|Z(u)|$. The opponent chooses the donor $t$, among those with $b(t)\ge1$.
\item \emph{New edge}: a held edge $tu$ becomes a new one-edge path, and $b(t),b(u)$ increase by one.
\end{itemize}
\item The play wins if $F$ is used up with $\sum b\le n$ and $b(r)\ge1$.
\end{enumerate}
\end{definition}
Every move is an instance of (B), or the addition of a new path. So a winning play is a certificate
for (R) at $(G,r)$ that never exhibits a path. Every Lov\'asz--Fan argument in the project has
this form: Lov\'asz from Fan, the two descent steps, the star lemma, Pyber's theorem, the
six-clique schedules and the $d=6$ frontier certificates~\cite{ES,D6}. The game wins every root
of every connected graph of order $6$. It loses on $50$ connected graphs of order $8$ with
$|F|\le4$. The experiment log lists five that still lose with $|F|\le 8$~\cite[E13]{EXP}:
\begin{center}
\texttt{GsRvnk}, \texttt{GsZvnk}, \texttt{Gqrvn[}, \texttt{Gqzvn[}, \texttt{Gqzl|\char123}.
\end{center}
The smallest, \texttt{GsRvnk}, consists of a clique $Q=\{q_1,\dots,q_4\}$ whose vertices have
degree $6$, a vertex $x$ adjacent to $Q$ and to a leaf $\ell$, and two false twins $s,s'$
adjacent exactly to $Q$. It loses at the roots $s,s'$. The same experiment showed that inductive
seeds, i.e.\ the rooted test on $G-F$ with an adversarial half-set, do not repair it.

These five graphs are \emph{not critical}: $p=3<n/2$. For them (R) is trivial, because a minimum
partition can be split at an interior visit to $r$. The game does not see this. In the game it is
the win condition $\sum b\le n-2$, which we add from now on (the \emph{split rule}). It never
applies to a critical one. It makes (R) trivial for every non-critical graph, but the game must
still reach $\sum b\le n-2$: with the split rule the five graphs above still lose at their losing
roots with $|F|\le8$ (audit A12). The critical losers of the original
game include the following.
\begin{itemize}
\item \texttt{Gqz\char126vw}, i.e.\ $K_8$ minus $\{03,12,23,34,56,67\}$. It is $4$-connected,
$p=4=n/2$, and the original game loses it at the root $3$ with $|F|\le8$ (a rerun of E13).
\item $K_8$ minus a perfect matching. It loses with $|F|\le8$ even with the move introduced
below (Theorem~\ref{thm:complete8}), and therefore also in the original game.
\end{itemize}
Among the $47$ graphs that still lose with the new move and $|F|\le4$, $19$ are critical
(\texttt{scripts/outputs/crit47.log}).

\subsection*{Results} Write \emph{path insertion} for the move
\begin{itemize}
\item \emph{New path}: a path $P$ all of whose edges are held is added as one new path, and $b$
increases by one at each end of $P$.
\end{itemize}
It is sound for the same reason as the new-edge move. The following are proved here or computed.
\begin{enumerate}
\item \emph{The hub lemma} (Lemma~\ref{lem:hub}): re-inserting four edges at an isolated vertex
costs one path, and the profile changes are controlled. \emph{Degree-four hubs}
(Theorem~\ref{thm:hub4}): if $d(s)=4$, every neighbour of $s$ is even, and at most one even vertex
lies outside $N[s]$, then (R) holds at $s$ and at every neighbour of $s$. A variant with one odd
neighbour (Theorem~\ref{thm:hub4odd}) closes the critical loser \texttt{Gqz\char126vw}. All are
formalised in Lean.
\item \emph{Completeness at order $8$} (Theorem~\ref{thm:complete8}, computation): the game with
path insertion wins every root of every connected graph of order $8$, with at most nine held
edges. The certificates found for the five losers all have the shape of the hub lemma.
\item \emph{P\'osa moves at the profile level fail} (Section~\ref{sec:posa}). We prove a rotation
lemma and \emph{Fan's alternative} (Lemma~\ref{lem:alt}). The alternative dominates rotation
as a game move (Section~\ref{sec:posa}). Neither repairs any losing root, with up to two uses per play and
$|F|\le4$, or one use and $|F|\le6$. We explain why: the losers' obstruction is budget
allocation, not exchange.
\item \emph{Degree six} (Section~\ref{sec:six}): at a hub of degree six the spare path pays for
three receipts of a pure star, and four are needed. The fourth sits at Fan's equality. Two hand
arguments settle it there. On the order-$10$ analogue of \texttt{GsRvnk} the receipt is forced by
an append-or-through argument and a count of edges (Proposition~\ref{prop:A62}). On $K_8$ minus a
perfect matching it always fails, and the rigid configuration is a flower of three Hamiltonian
cycles, which recombines (Proposition~\ref{prop:K8M}; a re-derivation of the case $m=4$
of~\cite[Thm.~10.21]{SP}). The game compensates instead by holding
more edges and routing the inserted path through the hub along them. A profile-level
receipt-or-flower move with an edge count (Lemma~\ref{lem:rfcount}), together with the symmetry
of the root's orbit, wins all these graphs with a single held star.
\item \emph{Arbitrary degree} (Section~\ref{sec:general}). A hub of degree $d$ whose neighbours
carry $\sum b\ge2d-5$ ends costs one path (Lemma~\ref{lem:genhub}). Relays from even
non-neighbours raise the supply (Theorem~\ref{thm:relayhub}). Two hubs with one common neighbour
share one new path (Lemma~\ref{lem:twohub}). The profile-level lemmas are formalised. The
hub-strategy conjecture (Conjecture~\ref{conj:hub}) holds at order $8$ with $c=2$. A sample of
order $10$ needs $c\ge3$, and the proved theorems cover none of its roots.
\end{enumerate}
For the programme this corrects connections~\S R. Lov\'asz seeds with Fan transfers were
incomplete at order $8$ because the game could not spend Lov\'asz's spare path on a path of length
two, not because of Fan's inequality. The next difficulty is the degree-six hub, which is the
``short by one'' of the $d=6$ frontier~\cite[Remark~3.5]{D6}. There the rigid configuration at
Fan's equality must be refuted or recombined with information outside the profile, or bypassed by
longer inserted paths.

\section{The spare path and the hub lemma}\label{sec:hub}

\subsection*{The spare path} Let $G$ have even order $n$, let $F\subseteq E(G)$, and suppose some
vertex $s$ is isolated in $K=G-F$. Then $K$ has at most $n-1$ vertices of positive degree. If
these include at most one even vertex, Lov\'asz's partition has $|\Od(K)|/2\le\lfloor
(n-1)/2\rfloor=n/2-1$ paths. In the game this means $\sum b\le n-2$ at the start. One new path,
of any length, can be afforded. Fan transfers never change the number of paths, so the spare path
can only be spent by a new path. Before this note the game allowed new paths of length one only.

\begin{lemma}[Hub lemma]\label{lem:hub}
Let $D$ be a partition of a graph $K$, and let $s$ be a vertex isolated in $K$. Let $a,c,d,x$ be
distinct vertices other than $s$, with $b_D(a),b_D(c),b_D(d)\ge1$. Then for some labelling
$\{y,y'\}=\{c,d\}$ the graph $K+sa+sc+sd+sx$ has a partition $D'$ with $|D'|=|D|+1$ and
\[
b_{D'}=b_D+2e_s-e_a-e_y+e_{y'}+e_x .
\]
In particular $b_{D'}(s)=2$ and $b_{D'}(x)=b_D(x)+1$.
\end{lemma}

\begin{proof}
\emph{Step 1.} $s$ has no neighbour in $K$, so $Z_D(s)=\emptyset$ and $b_D(a)\ge1>0$. A single
transfer of $sa$ towards $s$ gives a partition $D_1$ of $K+sa$ with $b_{D_1}=b_D+e_s-e_a$.

\emph{Step 2.} The only neighbour of $s$ in $K+sa$ is $a$, so $|Z_{D_1}(s)|\le1$, while
$b_{D_1}(c)+b_{D_1}(d)=b_D(c)+b_D(d)\ge2$. A grouped transfer with $T=\{c,d\}$ gives, for some
$y\in\{c,d\}$, a partition $D_2$ of $K+sa+sy$ with $b_{D_2}=b_{D_1}+e_s-e_y$.

\emph{Step 3.} Let $y'$ be the other vertex of $\{c,d\}$. The edges $sy'$ and $sx$ are not in
$K+sa+sy$, and $y'sx$ is a path. Adding it as a new path gives $D'=D_2+y'sx$, with $|D'|=|D|+1$
and the stated profile.
\end{proof}

In Lean this is \texttt{Ach.hub\_four}, stated for the end-vector bookkeeping \texttt{Ach} of the
insertion library. It holds for any partition $D$, not only a Lov\'asz partition. The vertex $x$ is
prescribed. The first donor $a$ is prescribed. Only the second donor is chosen by the adversary.

\begin{theorem}[Degree-four hubs]\label{thm:hub4}
Let $G$ have even order $n$ and let $s$ be a vertex of degree $4$ all of whose neighbours have
even degree. Suppose at most one vertex of $\Ev(G)$ lies outside $N[s]$. Then for every
neighbour $x$ of $s$, $G$ has a partition into at most $n/2$ paths with $b(s)=2$ and $b(x)=2$.
Consequently (R) holds at $s$ and at every neighbour of $s$.
\end{theorem}

\begin{proof}
Let $K=G-\{sw:w\in N(s)\}$, i.e.\ $G$ with $s$ isolated. A neighbour $w$ of $s$ has
$d_K(w)=d_G(w)-1$, which is odd. A vertex $w\notin N[s]$ has $d_K(w)=d_G(w)$. Hence the vertices
of positive even degree in $K$ are the even vertices of $G$ outside $N[s]$, and there is at most
one of them. By (A), $K$ has a partition $D$ with $b_D=d_K\bmod 2$. Every vertex of $K$ of odd
degree has one end, and $s$ has none. So $b_D(w)=1$ for each of the four neighbours of $s$, and
$2|D|=|\Od(K)|\le n-1$. Since $n$ is even and $|\Od(K)|$ is even, $2|D|\le n-2$.

Write $N(s)=\{a,c,d,x\}$ and apply Lemma~\ref{lem:hub}. The result is a partition $D'$ of $G$ with
$2|D'|=2|D|+2\le n$, $b_{D'}(s)=2$ and $b_{D'}(x)=b_D(x)+1=2$. For the last statement apply this
with any neighbour $x$ when $r=s$, and with $x=r$ when $r\in N(s)$.
\end{proof}

In Lean: \texttt{hub\_four\_ends}, and \texttt{hub\_four\_rooted} for the private-leaf form
$p(G+\text{leaf at }r)\le n/2$. The hypothesis on even vertices is stated as
$|\{v\in\Ev(G):v\ne s,\ sv\notin E\}|\le1$. Lov\'asz's theorem is the library's
\texttt{lovasz\_ends}, so the theorem is unconditional.

\begin{corollary}\label{cor:GsRvnk}
In \texttt{GsRvnk}, (R) holds at the roots $s$ and $s'$, where the original game loses.
\end{corollary}
\begin{proof}
$d(s)=4$ and $N(s)=Q$ consists of even vertices. The only even vertex outside $N[s]$ is $s'$.
Apply Theorem~\ref{thm:hub4} at $s$, and symmetrically at $s'$.
\end{proof}

Theorem~\ref{thm:hub4} has two features worth noting.
\begin{itemize}
\item Its hypothesis forces $|\Ev(G)|\le6$, so it lives in the regime of the open floor bound
$\mathrm{BP}_5$. It applies to vertices of degree four. In the apex class of the three-connected
deletion test $\delta\ge6$, so it does not apply there directly (Section~\ref{sec:six}).
\item Lemma~\ref{lem:hub} accepts any seed. Hence a floor bound for the odd-order graph $G-s$,
i.e.\ a partition into $n/2-1$ paths with an end at each neighbour of $s$, converts into the
rooted test at $s$ whenever $d(s)=4$. This is a small instance of the private-leaf duality
between the floor bound at odd order and (R) at even order.
\end{itemize}

\begin{remark}[Variants; the other four losers]\label{rem:variants}
The same scheme, receipts at an isolated hub followed by one new path through it, has several
variants. The certificates found for the other losers use the following (printed by
\texttt{scripts/certificate.py}).
\begin{itemize}
\item \texttt{Gqzl|\char123}, root $1$: the hub has degree $4$, but one neighbour is odd. One more
edge at that neighbour is held, so it is odd in $K$, and the new path continues along it:
$5\,1\,3\,6$ or $4\,1\,3\,6$.
\item \texttt{GsZvnk}, root $4$: a hub of degree $5$. It receives from $0$ and passes the end on
to $4$ by a single transfer. A grouped receipt and the new path $y'\,1\,7$ follow.
\item \texttt{Gqzvn[}, root $4$: the hub $3$ has degree $3$, and the new path $6\,3\,5\,7$
extends through a further held edge. The root receives by a single transfer from the seed.
\item \texttt{Gqrvn[}, root $2$: as \texttt{GsRvnk}.
\end{itemize}
A hub of degree $d$ with $b\ge1$ on its neighbours absorbs at least $\lceil d/2\rceil$ receipts
against an adversary. After $j$ receipts the supply is at least $d-j$ and at most $j$ neighbours
are blank. The remaining $d-\lceil d/2\rceil$ edges fit into one new path through the hub when
this number is at most $2$, i.e.\ $d\le5$. So the pure-star scheme closes hubs of degree at most
$5$, and at degree $6$ it is one receipt short (Section~\ref{sec:six}).
\end{remark}

\subsection*{A hub with an odd neighbour; the critical loser}
The variant used for \texttt{Gqzl|\char123} also closes the critical loser of the original game.

\begin{theorem}[Degree-four hubs with one odd neighbour]\label{thm:hub4odd}
Let $G$ have even order $n$, and let $s$ be a vertex of degree $4$ whose neighbours $a,c,d$ are
even and whose fourth neighbour $x$ is odd. Let $z\notin N[s]$ be an even neighbour of $x$, and
suppose at most one even vertex of $G$ lies outside $N[s]\cup\{z\}$. Then $G$ has a partition
into at most $n/2$ paths with $b(s)=2$ and $b(z)=2$. So (R) holds at $s$ and at $z$.
\end{theorem}
\begin{proof}
Let $K=G-\{sa,sc,sd,sx,xz\}$. In $K$ the vertices $a,c,d$ have odd degree $d_G-1$. The vertex
$x$ has odd degree $d_G(x)-2$, $z$ has odd degree $d_G(z)-1$, and $s$ is isolated. The vertices
of positive even degree in $K$ are the even vertices of $G$ outside $N[s]\cup\{z\}$, so there is
at most one. As in Theorem~\ref{thm:hub4}, (A) gives a partition $D$ of $K$ with $b=1$ at
$a,c,d,x,z$ and $2|D|\le n-2$. Steps 1 and 2 of Lemma~\ref{lem:hub} receive $sa$ and one of
$sc,sd$ at $s$. The remaining held edges $sy'$, $sx$, $xz$ form the path $y'\,s\,x\,z$, which is
added as a new path. Then $x$ is interior and keeps its end, $z$ gets a second end, $b(s)=2$, and
$|D'|=|D|+1\le n/2$.
\end{proof}

\begin{corollary}\label{cor:gqz}
(R) holds for \texttt{Gqz\char126vw} at the root $3$.
\end{corollary}
\begin{proof}
Take $s=3$, $N(3)=\{1,5,6,7\}$, $x=6$ (the odd neighbour) and $z=0$: then $06\in E$, $03\notin
E$, and $0$ is even. The even vertices are $0,1,3,4,5,7$, and the only one outside
$N[3]\cup\{0\}$ is $4$. As a play, $F=\{31,35,36,37,06\}$ and the Lov\'asz seed has $b=0$ at
$3,4$. The single transfer $1\to3$ comes first, then the grouped transfer from $\{5,7\}$ to $3$,
then the new path $0\,6\,3\,7$ or $0\,6\,3\,5$. The final profile is simple: it realises
$\{0,3,7\}$ or $\{0,3,5\}$ (\texttt{scripts/outputs/certificate\_Gqz.txt}).
\end{proof}
In Lean these are \texttt{Ach.hub\_four\_tail}, \texttt{hub\_four\_odd\_ends} and
\texttt{hub\_four\_odd\_rooted}. The instance is \texttt{gqz\_rooted\_three}, whose hypotheses
are checked by \texttt{decide} on $\mathrm{Fin}\,8$. The vertex $3$ has three even neighbours and
the graph is $4$-connected, so this is an instance of the branching base case
\texttt{BranchBase3} (with $k=3$) of the project's Lean reduction, by a simple profile.

\section{Completeness at order eight}\label{sec:eight}

\begin{theorem}[computation]\label{thm:complete8}
The Fan game with path insertion wins every root of every connected graph of order $8$ (all
$11{,}117$ graphs) with $|F|\le9$. Held sets with $|F|\le4$ suffice for all but $47$ graphs.
For those $47$ the least sufficient size, maximised over the roots, is as follows:
\begin{itemize}
\item $|F|=5$ for $35$ graphs;
\item $|F|=6$ for $6$ graphs;
\item $|F|=7$ for $4$ graphs;
\item $|F|=8$ for \texttt{Guv]\char126\char123}, at one root;
\item $|F|=9$ for $M_8=K_8$ minus a perfect matching, at every root. Every held set of size at most
$8$ loses.
\end{itemize}
\end{theorem}

The $47$ graphs are mostly the ones that need large held sets in any case, such as $K_{2,6}$. With
$|F|\le4$ the game with path insertion has $47$ losing graphs and the original game has $50$. The
three graphs repaired at $|F|\le4$ are \texttt{GsRvnk}, \texttt{Gqrvn[} and \texttt{GswNvs}. The
seven dense graphs \texttt{Gr\char126v\char126w}, \texttt{Gut\char126\char126w},
\texttt{Gu\char94v\char126w}, \texttt{Guv]\char125\char123}, \texttt{Guv]z\char123},
\texttt{Guv]\char126\char123} and \texttt{Gut\char126vs} (\texttt{outputs/dense7.g6}) were run one
root per automorphism orbit. Their complements have at most six edges. So do the complements of
three of the other $40$, \texttt{Gqz\char126vw}, \texttt{Gqzn\char126w} and
\texttt{Gqz\char94\char126w}, which were run at every root.
For $M_8$ the winning held set is the star at the root plus three edges. The play is a relay into
the hub, receipts, one donation from the hub, and the \emph{new path $1\,4\,2\,0\,7$ of length
four through the hub}, routed along the extra held edges (\texttt{outputs/certificate\_K8M.txt}).
Every certificate was found by \texttt{scripts/fan\_game.py}. Scripts and logs are listed in
Section~\ref{sec:scripts}.

Theorem~\ref{thm:complete8} was computed without the split rule. With it, (R) is trivial for the
$28$ non-critical graphs among the $47$. The game itself does not change: with the split rule it
still loses all $47$ with $|F|\le4$, and the least sizes above are the same (audit A12). The other
$19$ are critical, and for them the theorem says that path insertion alone suffices. Independently
of the game, \texttt{scripts/rootcheck.py} confirms (R) at every root of the five losers and of
$A_{6,2}$, $B$ and $B'$ below by a direct CP-SAT search for partitions, and \texttt{scripts/crit.py}
decides criticality.

\section{P\'osa moves at the profile level}\label{sec:posa}

Fan's lemma is iterated P\'osa rotation. Two profile-level versions of the iteration suggest
themselves. Both are sound, and both are useless on the losers.

For a path $P$ of a partition $D$ with an end $x$ and with $u\in V(P)$, let $\pi_P(x)$ be the
neighbour of $u$ on $P$ on the side of $x$, i.e.\ the vertex before $u$ when $P$ is read from $x$.

\begin{lemma}[Rotation]\label{lem:rot}
Let $D$ be a partition of $K$, $u\ne v$, $uv\notin E(K)$ and $b_D(v)\ge1$. Then one of the
following holds.
\begin{enumerate}[label=(\alph*)]
\item $K+uv$ has a partition $D'$ with $|D'|=|D|$ and $b_{D'}=b_D+e_u-e_v$.
\item There is a set $W\subseteq N_K(u)$ with $|W|=b_D(v)$ such that for every $w\in W$ the graph
$K+uv-uw$ has a partition $D_w$ with $|D_w|=|D|$ and $b_{D_w}=b_D-e_v+e_w$.
\end{enumerate}
\end{lemma}
\begin{proof}
If some path ending at $v$ avoids $u$, append $uv$ to it, which is (a). Otherwise each of the
$b_D(v)$ paths ending at $v$ contains $u$. These paths are distinct, since a path has distinct
ends. Their edges $u\,\pi_P(v)$ are distinct, so $W=\{\pi_P(v)\}$ has $b_D(v)$ elements, and $v\notin
W$ because $uv\notin E(K)$. Write $P=v\cdots w\,u\cdots$ with $w=\pi_P(v)$, and replace $P$ by
$w\cdots v\,u\cdots$. This is a path on the same vertex set, using $vu$ instead of $wu$, and it
moves the end from $v$ to $w$.
\end{proof}

Iterating rotations along Fan's chains gives more. If the transfer fails, then at least as many
\emph{blank} neighbours of $u$ as the supply are reachable, and each of them can be given an end.

\begin{lemma}[Fan's alternative]\label{lem:alt}
Let $D$ be a partition of $K$, let $u$ be a vertex, and let $T$ be a set of non-neighbours of $u$
with $u\notin T$. Then one of the following holds.
\begin{enumerate}[label=(\alph*)]
\item Some $t\in T$ with $b_D(t)\ge1$ gives an end to $u$: $K+tu$ has a partition $D'$ with
$|D'|=|D|$ and $b_{D'}=b_D+e_u-e_t$.
\item There is a set $W\subseteq Z_D(u)$ with $|W|\ge\sum_{t\in T}b_D(t)$ such that for each
$z\in W$ there are $t\in T$ with $b_D(t)\ge1$ and a partition $D_z$ of $K+tu-zu$ with
$|D_z|=|D|$ and $b_{D_z}=b_D-e_t+e_z$.
\end{enumerate}
\end{lemma}
In particular, if $\sum_Tb>|Z_D(u)|$, then (b) is impossible. This is (B).

\begin{proof}
A \emph{chain} is a sequence $t=x_0,x_1,\dots,x_k$ of distinct vertices with $t\in T$ such that for
each $i<k$ some path $P_i\in D$ has an end at $x_i$, contains $u$, and has
$\pi_{P_i}(x_i)=x_{i+1}$. Let $R$ be the set of last vertices of chains.

\emph{Executing a chain.} For a chain from $t$ to $z$ with $k\ge1$, the graph $K+tu-zu$ has a
partition with the same number of paths and profile $b_D-e_t+e_z$. Rotate $P_0$ at $u$ as in
Lemma~\ref{lem:rot}: the end moves from $x_0$ to $x_1$, and the edge $x_1u$ is replaced by $x_0u$.
The other paths are unchanged. In $P_0$ the side of $u$ towards its other end is unchanged, and
$x_1$ is interior to $P_0$, since it lies between $x_0$ and $u$. Hence $x_1,\dots,x_k$ is a chain
of the new partition: its steps use paths other than $P_0$, or $P_0$ read from its unchanged end
$\ne x_0$ (the $x_i$ are distinct). Induct on $k$. This is the reroute step of Fan's proof
(\texttt{PathDecomp.reroute\_move}, \texttt{fanChain\_transport} in the Lean library).

\emph{Case 1: some chain ends at a vertex $x_k$ with an end on a path avoiding $u$.} Execute the
chain and append $x_ku$ to that path. This gives (a) for $t=x_0$, and $b_D(t)\ge1$ because $P_0$
ends at $t$ (or $k=0$ and $t$ itself has the avoiding end).

\emph{Case 2: otherwise.} Every end at a vertex of $R$ lies on a path through $u$. Map each such
end $(P,x)$ to the edge $u\,\pi_P(x)$. Here $\pi_P(x)\in R$. Indeed, let $x_0,\dots,x_k=x$ be a
chain. If $\pi_P(x)$ is not among the $x_i$, it extends the chain. Otherwise it equals some $x_i$,
which is the last vertex of the chain $x_0,\dots,x_i$. The map is injective, since paths are edge-disjoint and the two ends
of one path lie on the two sides of $u$. Its image consists of edges from $u$ to $N(u)\cap R$. Also
$R\subseteq T\cup N(u)$, and $T\cap N(u)=\emptyset$. So
\[
\sum_{t\in T}b(t)+\sum_{w\in N(u)\cap R}b(w)\ \le\ \sum_{x\in R}b(x)\ \le\ |N(u)\cap R|.
\]
Every $w\in N(u)\cap R$ with $b(w)\ge1$ contributes at least $1$ on the left. Hence
$\sum_Tb\le|Z_D(u)\cap R|$. Put $W=Z_D(u)\cap R$. A vertex $z\in W$ is not in $T$, so it is the end
of a chain with $k\ge1$ from some $t\in T$. Its first step uses a path ending at $t$, so
$b_D(t)\ge1$. Executing the chain gives (b).
\end{proof}

Lemma~\ref{lem:alt} is the ``rigidity at equality'' asked for in~\cite[\S8]{D6}. When
$\sum_Tb=|Z(u)|$ and the transfer fails, every blank neighbour of $u$ is reachable, and each one can
be \emph{activated}: it is given an end, while the held edge moves to it.

\subsection*{As game moves} In the game, Lemma~\ref{lem:rot} gives the following move for a held
edge $uv$: the direct outcome must win, and the opponent may also choose $W\subseteq N_H(u)$ with
$|W|=b(v)$, after which we pick $w\in W$. Lemma~\ref{lem:alt} gives the same kind of move with
$W\subseteq Z_H(u)$, $|W|\ge\sum_Tb$, and any $T$. For $T=\{v\}$ the outcomes are the same states,
and the losing choices in $Z_H(u)$ are among those in $N_H(u)$. So the alternative wins whenever
rotation does: as moves, rotation is dominated. A rotated or activated state has the same number of
held edges, so plays are bounded by a budget of such moves.

\begin{proposition}[computation]\label{prop:posa}
On the five losers, adding rotations (at most two per play, $|F|\le4$) or Fan's alternative (at most
two per play with $|F|\le4$, or one per play with $|F|\le6$) to the original game repairs
\emph{none} of the losing roots: $s,s'$ in \texttt{GsRvnk}, and the single losing root of each of
the other four.
\end{proposition}
The moves are not idle. With $|F|\le4$ the original game loses \texttt{GsRvnk} at every root, and
with either move added it wins there at the six roots other than $s,s'$ (audit A12). These roots
are won by the original game with $|F|=6$.

\subsection*{Why} Every P\'osa-type move keeps the number of paths. The losers need the spare
path to be spent on a path of length two through the hub. In the original game this can only be
imitated in two steps. In \texttt{GsRvnk} the two receipts leave $b(s)=2$. First a new edge $cs$
is added, so $b(s)$ rises to $3$. Then the
transfer of $sd$ towards $d$ is needed. In the true partition this is an append: the path $cs$
ends at $s$ and avoids $d$. The profile has forgotten which of the three paths ending at $s$ is
$cs$, and both Fan's inequality and the rotation adversary now place $d$ on all of them. Path
insertion is exactly the memory of one explicit path. The missing ingredient at order $8$ is thus
budget allocation, not exchange.

\section{Degree six: the pure star is one receipt short}\label{sec:six}

The $d=6$ frontier of the three-connected test has $\delta\ge6$, and its residual case is a $K_6$
of even vertices with full exterior vertices~\cite[\S\S1, 3]{D6}. Remark~\ref{rem:variants} shows what
goes wrong. A hub of degree six with $b\ge1$ on its neighbours absorbs three receipts against an
adversary. Closing it with the spare path needs four receipts and one new path. The fourth receipt
has supply $3$ against three blank neighbours: it is exactly at Fan's equality.

Degree-six hubs are treated at the level of paths in~\cite[\S\S12.14, 13]{SP}. Let $y$ have degree
six and five even neighbours. Suppose (R) fails at a root near $y$ (Thm.~12.29) or at distance at
least three (Thm.~13.9). Then $G-y$ (or a slight modification of it) has more than $h$ paths and
misses at least $h+1$ edges, where $n=2h+2$. The proofs assign ending paths of a partition of
$G-y$ to the neighbours of $y$ and append the spokes. Doubly assigned paths close into cycles
through $y$, which a short path reopens, by the short-path-and-cycles lemma of Chu, Fan and Zhou.
These are the path-level forms of the receipts and flowers below. The profile game does not see
the assignments.

The game gets around this by holding further edges. Relays then feed the hub, and the new path is
routed through the hub along the extra held edges. For $M_8$ (Theorem~\ref{thm:complete8}) the
inserted path is $1\,4\,2\,0\,7$. We first give two hand arguments that close the pure star
itself at equality. Both use information that the profile does not record.

The order-$10$ analogue of \texttt{GsRvnk} is the graph $A_{6,2}$. It has a clique
$Q=\{q_1,\dots,q_6\}$ of vertices of degree $8$, a vertex $x$ adjacent to $Q$ and to a leaf
$\ell$, and false twins $s,s'$ adjacent exactly to $Q$. So $n=10$, $|E|=34$ and
$\Ev=Q\cup\{s,s'\}$.

\begin{lemma}[Append or through]\label{lem:through}
Let $D$ be a partition of $H$ and $t,u$ non-adjacent with $b_D(t)\ge1$. If no path ending at $t$
avoids $u$ then every path ending at $t$ contains $u$. If moreover $b_D(u)=d_H(u)$, then every
path containing $u$ ends at $u$, so every path ending at $t$ joins $t$ to $u$.
\end{lemma}
\begin{proof}
The first sentence restates the hypothesis. A path through $u$ that does not end at $u$ uses two
edges at $u$ that carry no end at $u$, so $b_D(u)\le d_H(u)-2$.
\end{proof}

\begin{proposition}\label{prop:A62}
$A_{6,2}$ has a partition into $5$ paths with $b(s)=4$ and $b=2$ at two core vertices. Hence (R)
holds at every vertex of $A_{6,2}$.
\end{proposition}
\begin{proof}
Let $K=G-s$, with $s$ isolated. Then $d_K(q)=7$, $d_K(x)=7$, $d_K(\ell)=1$, and $s'$ is the only
even vertex of positive degree. By (A), $K$ has a partition $D$ into $4$ paths with $b=1$ on
$Q\cup\{x,\ell\}$. Exactly as in Lemma~\ref{lem:hub}, a single transfer and two grouped transfers
towards $s$ (supplies $1>0$, $5>1$, $4>2$) produce a partition $D_3$ of $H=K+sy_1+sy_2+sy_3$ into
$4$ paths, for some distinct $y_i\in Q$. In $D_3$ we have $b(s)=3=d_H(s)$, $b=1$ on
$T=Q\setminus\{y_1,y_2,y_3\}$ and on $x,\ell$, and $b=0$ on $y_1,y_2,y_3,s'$.

\emph{Claim: some $t\in T$ has a path ending at $t$ that avoids $s$.} Suppose not. By
Lemma~\ref{lem:through} each of the three ends at $T$ lies on a path joining $T$ to $s$. These are
three distinct paths, since each has only one end in $T$, and so they are all three paths ending at
$s$. The fourth path has its ends at the remaining vertices with ends, $x$ and $\ell$. As $\ell$ is
a leaf at $x$, it is the edge $\ell x$. So the three $s$-paths cover $E(H)-\ell x$, which has
$34-3-1=30$ edges. But they avoid $\ell$, so each has at most $8$ edges, and $3\cdot8<30$. This
contradiction proves the claim.

Append $ts$ to that path. Now $b(s)=4$, and the two remaining edges $st'$, $st''$ at $s$ form the
new path $t'st''$. The result has $5=n/2$ paths, $b(s)=4$ and $b(t')=b(t'')=2$. Every permutation
of $Q$, fixing $x,\ell,s,s'$, is an automorphism of $A_{6,2}$, and so is the swap of $s$ and $s'$.
So (R) holds at every core vertex and at $s'$. For odd roots any partition into $5$ paths has an end
there by~\eqref{eq:parity}.
\end{proof}
(An earlier draft stated (R) only at $s,s'$ and the odd vertices, and listed the core roots as open.
The emc-four session pointed out that the symmetry of the core settles them~\cite[\S7]{ES}.)

The argument is a two-end phenomenon at equality, but it is not a transfer rule: it combines the
rigidity of Lemma~\ref{lem:through} with the number of edges, which the profile does not record.
In the game, the alternative of Lemma~\ref{lem:alt} at this state only activates a blank $y_i$. That
produces the same configuration with the labels permuted.

\begin{proposition}[computation]\label{prop:A62game}
The Fan game with path insertion, with $F$ a full star plus at most three further edges, loses
$A_{6,2}$ at every root, the odd ones included. (The CP-SAT check confirms (R) at every root.)
\end{proposition}

For the star plus three edges the solver ran one root per automorphism orbit, at a core vertex
and at $x$. A play that wins at any root wins at the odd root $x$, so a loss at $x$ means that no
play wins at all. At order $10$ a search over all held sets of size $\ge6$ is out of reach of the
solver, so $A_{6,2}$ is a loser only for hub-shaped held sets. It is exactly the case that
Proposition~\ref{prop:A62} settles by hand with the pure star. A sound count at Fan's equality,
together with the symmetry of the core, makes the game win it at every root
(Lemma~\ref{lem:rfcount} and Proposition~\ref{prop:rfgame} below). Three further test graphs of
order $10$ were run with the same held sets (\texttt{scripts/families.py}).
\begin{itemize}
\item $A_{4,4}$ ($Q=K_4$ with four false twins): it wins, and the twins need a star plus two
further held edges.
\item $B$ (a $K_6$ core $U$ of even vertices, two adjacent full vertices, and two vertices
each adjacent to half of $U$): it wins with a star plus one edge. The half vertices have small
contact, which makes them resources in the sense of~\cite[\S4]{D6}.
\item $B'$ (a $K_6$ core, a full vertex $x_1$ with a leaf, and two adjacent full vertices): this
is the residual case Q2 of the $d=6$ frontier, all exterior vertices being full. It loses at every
root with a star plus at most two edges. With a star plus three edges it wins at the core roots
and at $x_1$. The remaining roots are odd, and a play that wins at one root wins at every odd
root, by~\eqref{eq:parity}. So $B'$ wins at every root (\texttt{outputs/famhub\_B2\_3.log}). The
certificate for a core root holds the star at $x_1$ and a core triangle. Its play is a relay
through the hub, then receipts, then the new path $0\,2\,1\,6\,5$ through the hub
(\texttt{outputs/certificate\_B2.txt}). The count of Proposition~\ref{prop:A62} would not have
refuted the rigid case here: at the fourth receipt the fourth path joins the two other full
vertices, and $31$ edges meet a capacity of $3\cdot8+7=31$. The emc-four session closes $B'$ by
hand at every root from exactly this equality~\cite[Prop.~68]{ES}. Since every bound is attained,
the cycles through the hub are Hamiltonian, and its receipt-or-flower theorem~\cite[Thm.~67]{ES}
recombines them.
\end{itemize}

The rigid configuration can also be useful in itself. Here is the order-$8$ graph on which the
hub argument meets it head-on.

\begin{lemma}[Flowers]\label{lem:flower}
Let $C_1,\dots,C_j$ be edge-disjoint cycles, and let $e_i\in E(C_i)$ be edges that are the
consecutive edges of a path $L$. Then $C_1\cup\dots\cup C_j$ has a partition into $j+1$ paths with
two ends at every vertex of $L$ and no end elsewhere.
\end{lemma}
\begin{proof}
The paths are $C_i-e_i$ for $i\le j$, and $L$. The path $C_i-e_i$ ends at the two ends of $e_i$.
An interior vertex of $L$ lies on two of the $e_i$ and is interior to $L$. An end of $L$ lies on
one $e_i$ and ends $L$.
\end{proof}
Lemma~\ref{lem:flower} is elementary and not new. The cycle-opening results of the archive go
much further for the conclusion that matters for (R), namely two ends at the common vertex.
\begin{itemize}
\item Any $j\le7$ edge-disjoint cycles through a vertex $x$ partition into at most $j+1$ paths,
two of them ending at $x$~\cite[Thm.~38]{MR}. This extends Chu, Fan and Zhou's $j+1$ count.
\item The same holds for twelve nearly spanning cycles~\cite[Thm.~1]{CC}.
\item Larger families are treated under intersection conditions~\cite[Thm.~13.18]{SP}.
\end{itemize}
The Hamiltonian flowers below are special cases. What our version adds is only the prescribed
ends at the vertices of a transversal path.

\begin{proposition}\label{prop:K8M}
Let $M_8$ be $K_8$ minus a perfect matching. Then (R) holds at every vertex of $M_8$, by
Lov\'asz's partition of $M_8-s$, three Fan receipts at $s$, and Lemma~\ref{lem:flower}.
\end{proposition}
\begin{proof}
$M_8$ is $6$-regular and vertex-transitive, so let $s$ be the root and $s^*$ its non-neighbour. In
$K=M_8-s$ the six neighbours of $s$ have degree $5$ and $s^*$ is the only even vertex. By (A), $K$
has a partition into $3$ paths with $b=1$ on $N(s)$. As in Lemma~\ref{lem:hub}, three receipts at
$s$ (supplies $1>0$, $5>1$, $4>2$) give a partition $D_3$ of $H=K+sy_1+sy_2+sy_3$ into $3$ paths,
with $b(s)=3$ and $b=1$ on the three other neighbours $T$. The three paths have six ends, three at
$s$ and three in $T$. A path has at most one end at $s$, so every path joins $s$ to $T$. The held
edges $st$, $t\in T$, close them into three edge-disjoint cycles $C_1,C_2,C_3$ through $s$ that
cover $E(M_8)$. A cycle has at most $8$ edges and $|E(M_8)|=24$, so the $C_i$ are Hamiltonian.

Choose a neighbour $a$ of $s$ on $C_1$, and put $e_1=sa$. The cycle $C_2$ passes through $a$. Its
two edges there avoid $s$, because $sa\in C_1$. Choose one, $e_2=aw$. The cycle $C_3$ passes
through $w$. Its two edges there avoid $a$, because $aw\in C_2$, and at most one of them goes to
$s$. Choose $e_3=wz$ with $z\ne s$. Then $s,a,w,z$ are distinct and $L=s\,a\,w\,z$ is a path.
Lemma~\ref{lem:flower} gives $4=n/2$ paths with $b(s)=2$.
\end{proof}

Proposition~\ref{prop:K8M} is a re-derivation. (R) at every root of every cocktail party graph
$\mathrm{CP}(m)=K_{2m}$ minus a perfect matching, $m\ge4$, is~\cite[Thm.~10.21]{SP}. Its proof
opens the $m-1$ Walecki Hamilton cycles, each detoured through the mate of the root, along the
path $Q=(t,1,\dots,m-1)$. This is a flower with a transversal path, for every $m$. The same theorem
shows that Fan's numerical test fails in every partition within the budget, after any edge is
deleted. The Lean theorem \texttt{cocktail\_fan\_barrier} is a formal counterpart of that half.
New here is only that the flower arises from a Lov\'asz seed, three receipts and an edge count,
i.e.\ inside the game.

With the pure star the fourth receipt \emph{always} fails: all three paths end at the hub, so Fan's
chain from $t\in T$ stops at a blank donor. The spare path is spent on recombining the flower. (The
game, which cannot see the flower, needs nine held edges.) This is a two-end move at equality. It moves the ends from $T$ to the transversal $L$, and it is again
invisible in the profile: it needs the three cycles, whose existence follows from a count of
edges.

\subsection*{A profile-level flower move}
In Proposition~\ref{prop:K8M} everything needed for the flower can be read off $(H,b)$. It needs
only a transversal lemma for Hamiltonian cycles.

\begin{lemma}[Transversals]\label{lem:transversal}
Let $C_1,\dots,C_j$, $j\le4$, be edge-disjoint Hamiltonian cycles on a common vertex set, and let
$u$ be a vertex. Then there are distinct vertices $u=a_0,a_1,\dots,a_j$ such that the edges
$a_{i-1}a_i$ lie in distinct cycles.
\end{lemma}
\begin{proof}
Let $N$ be the number of vertices. Since $u$ has degree $2j$ in the union, $N\ge2j+1$. Choose
$a_1,a_2,\dots$ in turn, each time along a cycle not used before. From $a_{i-1}$ the
$j-i+1$ unused cycles offer $2(j-i+1)$ neighbours, and these are distinct because the edges at
$a_{i-1}$ are distinct. None of them is $a_{i-2}$, since the edge $a_{i-2}a_{i-1}$ lies in a used
cycle. So at most $i-2$ of them, namely $a_0,\dots,a_{i-3}$, are forbidden. For $j\le3$ we have
$2(j-i+1)>i-2$ at every step. For $j=4$ the same holds at steps $1$ to $3$. At step $3$ there
are four candidates, of which only $a_0$ may be forbidden, so there are at least three valid
choices of a cycle and a vertex. At step $4$ only the last cycle $C$ is left. The choice fails
exactly when the two neighbours of $a_3$ on $C$ are $a_0$ and $a_1$, i.e.\ when $a_3$ is a common
neighbour of $a_0$ and $a_1$ on $C$. On a cycle of length $N\ge5$ two vertices have at most one
common neighbour. A step-$3$ choice along the cycle $C_A$ leaves $C_B$ last, and a choice along
$C_B$ leaves $C_A$ last. So at most one valid choice of each kind fails, and at least one of the
at least three valid choices succeeds.
\end{proof}

\begin{lemma}[Flower move]\label{lem:flowermove}
Let $D$ be a partition of $H$, let $u$ be a vertex, and let $T\not\ni u$ be a set of $j\le4$
non-neighbours of $u$. Suppose that $b_D(u)=d_H(u)=j$, that $b_D=1$ on $T$, that
$\sum_vb_D(v)=2j$, and that $|E(H)|+j=j\,|V'|$, where $V'$ is the set of vertices of positive
degree in $H+uT$. Then for some path $L=u\,a_1\cdots a_j$ of $H+uT$, the graph $H+uT$ has a
partition into $j+1$ paths with $b=2$ on $V(L)$ and $b=0$ elsewhere.
\end{lemma}
\begin{proof}
$D$ has $j$ paths. Each has at most one end at $u$, and $u$ has $j$ ends, so each path has one end
at $u$. Its other end lies in $T$, the only other vertices with ends, and since $b=1$ on $T$ the
$j$ paths end at distinct $t\in T$. Adding $ut$ to the path ending at $t$ gives $j$
edge-disjoint cycles through $u$ that cover $H+uT$. Each has at most $|V'|$ edges, and together
they have $j|V'|$, so each is Hamiltonian on $V'$. Lemmas~\ref{lem:transversal}
and~\ref{lem:flower} finish the proof.
\end{proof}

As a game move (option \texttt{flower} of \texttt{fan\_game.py}), the hypotheses are checked on
$(H,b)$. The opponent picks $L$ among all paths of length $j$ from $u$ in $H+uT$, and every
outcome must win. The move inserts $j$ held edges at once. This is what the barrier theorems of
the emc-four session (\texttt{FanBarrier.lean}) require for the cocktail party graphs
$\mathrm{CP}(m)=K_{2m}$ minus a perfect matching. \texttt{cocktail\_fan\_barrier} shows that for
$m\ge4$ no certificate whose last step re-inserts one edge, by a Fan transfer or as a one-edge path,
works. \texttt{cocktail\_insertion\_barrier} shows that if the last step adds one new path made of
held edges to a partition of $\mathrm{CP}(m)-F$ into at most $m-1$ paths, then $|F|\ge m-1$. The
flower's final partition is $\{C_i-e_i\}\cup\{L\}$, so it is of this form with $F=E(L)$ and
$|F|=m-1$. It attains the bound.

\begin{proposition}[computation]\label{prop:flowergame}
With the flower move, the game with path insertion wins $\mathrm{CP}(4)$ and $\mathrm{CP}(5)$ at
every root with $F$ a single star, i.e.\ $|F|=2m-2$. Without it, $\mathrm{CP}(4)$ needs nine held
edges and $\mathrm{CP}(5)$ loses with the star. Elsewhere in the test bed the move adds nothing.
$A_{6,2}$ and $B'$ still lose with a star plus at most one edge. Five of the seven dense order-$8$
graphs still lose with $|F|\le6$. Five randomly sampled critical graphs of order $10$ with a vertex
of degree $6$ win at the same held-set size with or without it, with a star plus at most two
edges.
\end{proposition}
For $\mathrm{CP}(m)$ the argument is uniform up to the transversal. Lov\'asz on $\mathrm{CP}(m)-s$
gives $m-1$ paths with $b=1$ on $N(s)$. Then $m-1$ receipts follow (supplies $1>0$, $2m-3>1$,
$\dots$), and the flower with $j=m-1$ Hamiltonian cycles finishes. Lemma~\ref{lem:transversal}
covers $m\le5$, and a transversal lemma for every $j$ would give the flower move in the game for
every $\mathrm{CP}(m)$, each vertex of which is an instance of \texttt{BranchBase3} with $k=m$.
(R) itself holds for every $\mathrm{CP}(m)$ by~\cite[Thm.~10.21]{SP}.

\subsection*{Receipt or flower, with a count}
Proposition~\ref{prop:A62}, Proposition~\ref{prop:K8M} and the closure of $B'$ are three outcomes
of one alternative: the receipt-or-flower theorem of~\cite[Thm.~67]{ES}. If the receipt fails at a
saturated hub, then the paths at the hub join $T$ to the hub, and the other paths are left
unchanged. The following edge count, which uses only $H$ and $b$, decides between the cases often
enough for the game.

\begin{lemma}[Receipt or flower, with a count]\label{lem:rfcount}
Let $D$ be a partition of $H$, let $u$ be a vertex with $b_D(u)=d_H(u)=j$, and let $T\not\ni u$ be a set
of $j$ non-neighbours of $u$ with $b_D=1$ on $T$. Let $V^+$ be the set of vertices of positive degree
in $H$, and let $\Lambda$ be the set of vertices of degree $1$ in $H$ outside $T\cup\{u\}$. Put
$\lambda=|V^+|-|\Lambda|-1$ and $m'=|D|-j$. Let $\mu=0$ if $m'=0$. If $m'=1$ and the two remaining
ends are a leaf of $H$ and its neighbour, let $\mu=1$. Otherwise let $\mu=m'(|V^+|-2)$.
\begin{enumerate}[label=(\alph*)]
\item If $|E(H)|>j\lambda+\mu$, then some $t\in T$ ends a path of $D$ that avoids $u$. Appending $tu$
gives a partition of $H+tu$ into $|D|$ paths with profile $b_D+e_u-e_t$.
\item If $|E(H)|=j\lambda+\mu$ and $j\le4$, then this holds, or the following does: for some path
$L=u\,a_1\cdots a_j$ of $H+uT$ with all vertices in $V^+\setminus\Lambda$, the graph $H+uT$ has a
partition into $|D|+1$ paths with profile $b_D-j\,e_u-\mathbf 1_T+2\cdot\mathbf 1_{V(L)}$.
\end{enumerate}
\end{lemma}
\begin{proof}
By~\cite[Thm.~67]{ES} (or Lemma~\ref{lem:through}), either the conclusion of~(a) holds, or the $j$
paths at $u$ join $T$ to $u$ and the other $m'$ paths avoid $u$. Suppose the second case holds. A
vertex of degree $1$ can only be an end, so a path joining $T$ to $u$ avoids $\Lambda$ and has at
most $\lambda$ edges. Each other path has at most $|V^+|-2$ edges. If $m'=1$ and its ends are a leaf
$\ell$ and the neighbour of $\ell$, it is the edge between them. So $|E(H)|\le j\lambda+\mu$. In~(a)
the second case is therefore impossible. In~(b) every bound is attained, so each path joining $T$
to $u$ is a Hamiltonian path of $V^+\setminus\Lambda$, and the cycles $C_t$ are Hamiltonian there.
Lemma~\ref{lem:transversal} gives a transversal from $u$, and Lemma~\ref{lem:flower} applied to
the cycles, with the other paths unchanged, gives the partition.
\end{proof}

For $A_{6,2}$ at its fourth receipt, $j=3$, $|V^+|=10$, $\Lambda=\{\ell\}$, $\lambda=8$, $m'=1$ with
ends $\ell,x$, and $\mu=1$. Then $31>25$, so case~(a) holds: this is Proposition~\ref{prop:A62}.
For $B'$, $\ell$ is isolated while $x_1\ell$ is held, so $|V^+|=9$, $\Lambda=\emptyset$,
$\lambda=8$, $m'=1$ and $\mu=7$. Then $31=24+7$, so case~(b) holds, as in~\cite[Prop.~68]{ES}. For
$\mathrm{CP}(m)$, $m'=0$, and~(b) is Lemma~\ref{lem:flowermove}.

As a game move (option \texttt{rf}), in case~(a) the opponent picks $t$. In case~(b) the opponent
picks either such a $t$ (when $m'\ge1$) or a path $L$. A second sound rule, option \texttt{orbit},
uses that (R) is invariant under automorphisms: a play may end with an end at any vertex of the
automorphism orbit of the root.

\begin{proposition}[computation]\label{prop:rfgame}
With path insertion, receipt or flower, and the orbit rule, the game wins $A_{6,2}$, $B'$,
$\mathrm{CP}(4)$ and $\mathrm{CP}(5)$ at every root with $F$ a single star. Without the orbit rule
it still wins every root except the core roots of $A_{6,2}$ and $B'$. At those the opponent can
make the root a donor, and the hand proofs also use the symmetry of the core. On ten randomly
sampled critical graphs of order $10$ with a vertex of degree $6$, the move changes nothing. Nine
win with a star plus one or two held edges, with or without it. The tenth is
\texttt{Inv\char126|\char126zbw}, i.e.\ $K_{10}$ minus $\{02,17,19,24,29,39,58\}$. It is
$6$-connected with eight even vertices, and with a star plus at most two held edges it loses at
the roots $2$ and $3$, with or without the move (\texttt{scripts/outputs/rf\_sample.log}). With a
star plus three edges it wins at both roots, again with or without the move.
\end{proposition}

The root $2$ of \texttt{Inv\char126|\char126zbw} is a hub of degree $6$ with two odd neighbours,
$1$ and $6$. Isolating it leaves five even vertices. Holding the edges $10$ and $64$ as well leaves
only $9$, so a Lov\'asz seed exists. At the fourth receipt, however, $|E(H)|=33$ while
$j\lambda+\mu=3\cdot9+8=35$. So Lemma~\ref{lem:rfcount} decides neither case, and the rigid
configuration is not refuted. The game instead holds the star at $2$ and the edges $01$, $34$,
$36$ ($|F|=9$). Its play is a relay $0\to1\to2$ into the hub, receipts at $2$, a second hub at
$3$, and a new path through both hubs, $4\,3\,6\,2\,8$ or $4\,3\,6\,2\,7$. When the second
hub's donor is $4$, the new path is $2\,6\,3$ or $3\,6\,2\,7$ and single transfers finish
(\texttt{scripts/outputs/certificate\_miss1.txt}). Every leaf of this certificate has ends at both
$2$ and $3$, so it also wins at the root $3$. So in this sample the miss is again a matter of
held-set size: a second hub along the inserted path does what the count cannot.

\texttt{Inv\char126|\char126zbw} is dense ($\delta=0.6n$), and it lies in a regime that the archive
treats at the level of paths. Its vertex $9$ has degree $6$, and $6$ is its only odd neighbour.
Deleting $9$ leaves $K_9$ minus $\{02,17,24,58\}$, with $h=4$ missing edges. By the rainbow
Hamiltonian decomposition theorem~\cite{JM} (repaired in~\cite{JMR}) this graph has a partition into
$4$ Hamilton paths. So~\cite[Thm.~12.29]{SP} with $y=9$ gives (R) at the roots $9$ (distance $0$,
$s=6$) and $6$ (distance $1$, $s=t=6$). At the roots $1,2,3$ (distance two, through $6$) the graph
$W$ of that theorem misses $h+1=5$ edges, and the theorem does not decide. The roots $0,4,5,7,8$ are
not of its form, because a neighbour other than $6$ would have to be the exceptional site. So the
large held sets this graph needs (here, and in Proposition~\ref{prop:hubtest}(c)) measure how the
end budgets of the profile game starve at high density, as for $\mathrm{CP}(4)$. They do not point
to a sparse obstruction, and we did not raise the held-edge caps further.

\begin{question}\label{q:count}
\begin{enumerate}[label=(\alph*)]
\item Is the game with path insertion, receipt or flower, and the orbit rule complete at order
$10$? \texttt{Inv\char126|\char126zbw} is the test case.
\item Lemma~\ref{lem:rfcount} bounds the other paths crudely. Is there a sharper profile-level
bound, for instance one that uses the degrees of the remaining ends, that decides the rigid case
more often?
\item Does Lemma~\ref{lem:transversal} hold for every $j$ (with $N\ge2j+1$)? That would give a
profile-level flower move for every $j$. The rooted test for every $\mathrm{CP}(m)$ is already
known~\cite[Thm.~10.21]{SP}.
\end{enumerate}
\end{question}

\section{Hubs of arbitrary degree}\label{sec:general}

The certificates of Sections~\ref{sec:hub}--\ref{sec:six} have one shape. Hold the star of a hub
$s$ and a few further edges. Take a Lov\'asz seed of the rest, which has a spare path because $s$
is isolated. Let $s$ receive, and spend the spare path on a new path through $s$. This section
proves that shape for every degree, in three forms (one hub, relays into the hub, two hubs), and
states the conjecture that it always suffices with a bounded number of further held edges.
Path-level versions of the shape are older. In~\cite[Cor.~39]{MR}, a new vertex $x$ is joined to
$d$ ends of a partition into $s$ paths. If $q$ of the spokes cannot be matched to distinct paths,
then $s+\lceil q/7\rceil$ paths suffice, with at least $d-2q+2\lceil q/7\rceil$ ends at $x$. The proofs
of~\cite[Thms.~12.29, 13.9]{SP} and the hub phases of~\cite{D6} have the same shape. What is new
below is the profile-level form, which the game can use, with Lov\'asz's spare path as the
payment.

\begin{lemma}[Receipts]\label{lem:receipts}
Let $D$ be a partition of $H$, let $s$ be a vertex, and let $T$ be a set of non-neighbours of $s$
with $s\notin T$ and $1\le b_D\le2$ on $T$. If $r\le|T|$ and $|Z_D(s)|+2r\le\sum_Tb_D+1$, then for
some $R\subseteq T$ with $|R|=r$ the graph $H+sR$ has a partition into $|D|$ paths with profile
$b_D+r\,e_s-\mathbf 1_R$.
\end{lemma}
\begin{proof}
Induction on $r$. Let $r\ge1$. Then $\sum_Tb\ge|Z(s)|+2r-1>|Z(s)|$, so by~(B) some $t\in T$ with
$b(t)\ge1$ gives an end to $s$. The only new neighbour of $s$ is $t$, and it is blank exactly when
$b(t)=1$. No other vertex of $N(s)$ changes its value. So the new partition has
$|Z'(s)|\le|Z(s)|+[b(t)=1]$, and $\sum_{T-t}b=\sum_Tb-b(t)$. Since $b(t)+[b(t)=1]=2$, the hypothesis
holds for $T-t$ and $r-1$.
\end{proof}
So each receipt lowers $\sum_Tb-|Z(s)|$ by at most two. For $b=1$ on $T$ this is the receipt
schedule of~\cite[Remark~3.5]{D6}.

\begin{lemma}[General hub lemma]\label{lem:genhub}
Let $D$ be a partition of $H$, let $s$ be isolated in $H$, and let $N$ be a set of $d\ge2$
vertices, $s\notin N$, with $1\le b_D\le2$ on $N$ and $\sum_Nb_D\ge2d-5$. Then $H+sN$ has a
partition into $|D|+1$ paths with profile $b_D+(d-2)e_s-\mathbf 1_R+e_y+e_{y'}$, where
$N=R\sqcup\{y,y'\}$.
\end{lemma}
\begin{proof}
Apply Lemma~\ref{lem:receipts} with $T=N$ and $r=d-2$. Here $Z(s)=\emptyset$, and
$2(d-2)\le\sum_Nb+1$. The two remaining edges $sy,sy'$ form the new path $y\,s\,y'$.
\end{proof}
If $d\le5$, the supply condition holds as soon as $b\ge1$ on $N$. At $d=6$ one neighbour needs
two ends. A pure-star Lov\'asz seed has $b=1$ on $N$, so it is short by exactly one: this is the
missing receipt of Section~\ref{sec:six}. For $d=4$, Lemma~\ref{lem:hub} is stronger, because
there the fourth neighbour needs no end.

A held edge $xz$ from an odd neighbour $x$ of the hub to an even vertex $z$ fixes the parity of
both, and it can be returned first, so that $x$ has two ends when the hub receives.

\begin{lemma}[Relays]\label{lem:relays}
Let $D$ be a partition of $H$, let $X$ be a set of vertices, and let $z$ be an injective map from
$X$ to vertices outside $X$ with $b_D(z(x))=1$ and $z(x)x\notin E(H)$ for $x\in X$. Suppose that
no $x\in X$ is adjacent in $H$ to a blank vertex or to a vertex of $z(X)$. Then
$H+\{z(x)x:x\in X\}$ has a partition into $|D|$ paths with profile
$b_D+\mathbf 1_X-\mathbf 1_{z(X)}$.
\end{lemma}
\begin{proof}
Apply the single transfers $z(x)\to x$ in any order. Earlier transfers raise $b$ on $X$ and lower
it on $z(X)$, and $x$ sees no vertex of $z(X)$. The edges added before also avoid $x$. So $x$ has
no blank neighbour, and its supply $b(z(x))=1$ exceeds $|Z(x)|=0$.
\end{proof}
These are forward relays of length one~\cite[Lemma~2.3(a)]{D6}, done side by side.

\begin{theorem}[Relay hub]\label{thm:relayhub}
Let $G$ have even order $n$, and let $s$ be a vertex of degree $d\ge3$. Let $x_1,\dots,x_q$ be the
odd neighbours of $s$, and let $z_1,\dots,z_q$ be distinct even vertices outside $N[s]$ with
$x_iz_i\in E(G)$. Put $F=\mathrm{star}(s)\cup\{x_iz_i:i\le q\}$. Suppose the following hold.
\begin{enumerate}[label=(\roman*)]
\item $G-F-s$ has at most one vertex $w$ of even degree (degree $0$ counts as even).
\item No $x_i$ is adjacent to $w$, or to a $z_j$ with $j\ne i$.
\item $q\ge d-5$.
\end{enumerate}
Then $G$ has a partition into $n/2$ paths in which $s$ is an end of $d-2$ paths. In particular
(R) holds at $s$.
\end{theorem}
\begin{proof}
$G-F-s$ has odd order $n-1$ and an even number of odd vertices, so it has exactly one even vertex
$w$. By~(A) it has a partition into $(n-2)/2=n/2-1$ paths with $b=1$ on the odd vertices and
$b(w)=0$. Regard it as a partition of $H=G-F$, in which $s$ is isolated. The even neighbours of
$s$, the $x_i$ and the $z_i$ are odd in $H-s$, so $b=1$ there. In $H$ the blank vertices are
$s$ and $w$, and $x_is\in F$. So by~(ii) Lemma~\ref{lem:relays} applies with $X=\{x_i\}$ and
$z(x_i)=z_i$. Afterwards $b=2$ on $X$ and $b=1$ on the other neighbours of $s$. So
$\sum_{N(s)}b=d+q\ge2d-5$ by~(iii). Now Lemma~\ref{lem:genhub} inserts the star at the cost of
one path, and $s$ gets $d-2\ge1$ ends.
\end{proof}

\begin{lemma}[Two hubs]\label{lem:twohub}
Let $D$ be a partition of $H$, let $s_1\ne s_2$ be isolated in $H$, let $N_1,N_2$ be disjoint sets
avoiding $s_1,s_2$, and let $m\notin N_1\cup N_2\cup\{s_1,s_2\}$. Suppose $1\le b_D\le2$ on
$N_1\cup N_2$, $|N_i|\ge1$ and $\sum_{N_i}b_D\ge2|N_i|-3$ for $i=1,2$. Then the graph
$H+s_1(N_1\cup\{m\})+s_2(N_2\cup\{m\})$ has a partition into $|D|+1$ paths with profile
$b_D+\sum_{i=1,2}\bigl((|N_i|-1)e_{s_i}-\mathbf 1_{R_i}+e_{y_i}\bigr)$, where
$N_i=R_i\sqcup\{y_i\}$.
\end{lemma}
\begin{proof}
Lemma~\ref{lem:receipts} at $s_1$ with $T=N_1$ and $r=|N_1|-1$ changes $b$ only at $s_1$ and on
$N_1$, and keeps $s_2$ isolated. Then the same at $s_2$ with $T=N_2$. The four remaining edges
form the new path $y_1\,s_1\,m\,s_2\,y_2$. Its vertices are distinct by the disjointness
hypotheses.
\end{proof}

\begin{corollary}\label{cor:twohub}
Let $G$ have even order $n$, and let $s_1,s_2$ be non-adjacent vertices with $2\le d(s_i)\le4$ and
exactly one common neighbour. Suppose that every vertex of $G-s_1-s_2$ has odd degree. Then $G$
has a partition into $n/2$ paths with $d(s_i)-2$ ends at $s_i$ ($i=1,2$). In particular (R) holds
at $s_i$ if $d(s_i)\ge3$.
\end{corollary}
\begin{proof}
By~(A), $G-s_1-s_2$ has a partition into $(n-2)/2$ paths with $b=1$ everywhere. Apply
Lemma~\ref{lem:twohub} with $m$ the common neighbour and $N_i=N(s_i)-m$. The supply condition
$|N_i|\ge2|N_i|-3$ holds because $|N_i|\le3$.
\end{proof}
Two hubs cost one path, the same as one. So a second hub is a way to repair parity: it removes the
even vertices that the first hub leaves among the neighbours of the second.

In Lean, Lemmas~\ref{lem:receipts}, \ref{lem:genhub}, \ref{lem:relays} and~\ref{lem:twohub} are
\texttt{Ach.hub\_receipts}, \texttt{Ach.hub\_general}, \texttt{Ach.relays} and
\texttt{Ach.two\_hub}. The step of Theorem~\ref{thm:relayhub} after the seed is
\texttt{Ach.relay\_hub}. The wrappers with Lov\'asz's theorem, which \texttt{hub\_four\_ends} provides
at degree four, are not formalised for Theorem~\ref{thm:relayhub} and Corollary~\ref{cor:twohub}.

\subsection*{The hub-strategy conjecture}
With $F=\mathrm{star}(y)\cup X$ the seed is always Lov\'asz's, since $y$ is isolated. So $G-F$ may
have at most one even vertex of positive degree. In $G-\mathrm{star}(y)$ the even vertices of
positive degree other than $y$ are the odd neighbours of $y$ that are not leaves, and the even
vertices outside $N[y]$. Let $e(y)$ be their number. An edge of $X$ changes the parity of its two
ends, so $|X|\ge\mathrm{lb}(y)=\lfloor e(y)/2\rfloor$.

\begin{conjecture}[Hub strategy]\label{conj:hub}
Let $G$ be a connected graph of even order, and let $y$ be an even vertex with at least three
even neighbours and no bridge $yt$ with $d(t)\ge2$ (the hypotheses on the root in
\texttt{BranchBase3}). Then the game with path insertion, receipt or flower, and the orbit rule
wins at $y$, with the simple target, for $F=\mathrm{star}(y)\cup X$ and some set $X$ of at most
$\mathrm{lb}(y)+c$ edges, where $c$ is an absolute constant.
\end{conjecture}
The point is that $c$ does not depend on $G$. Without a bound the statement says no more than
(R), since with $F=E(G)$ path insertion can insert any partition. At order $8$ the conjecture holds with $c=2$,
and $c=2$ is needed. At order $10$ our sample needs $c\ge3$
(Proposition~\ref{prop:hubtest}). So the form with $c=2$, which we first conjectured from the
order-$8$ data, is false.
The \emph{simple target} (option \texttt{simple}) is the profile of a realisable set containing
$y$: $b=1$ on odd vertices, $b\in\{0,2\}$ on even ones, $b(y)=2$, or $b(y)=0$ when the split rule
applies. This is what \texttt{Desc} and \texttt{BranchBase3} ask for. The conclusion of
Lemma~\ref{lem:genhub} is not simple when $d\ge6$, since then $b(s)=d-2\ge4$. A play that passes
through it and wins with the simple target must afterwards move ends away from the hub
(donations, in the language of~\cite{D6}). In
\texttt{Inv\char126|\char126zbw} (Proposition~\ref{prop:rfgame}) the second hub $3$ is the vertex
$N$ of a path $O\!-\!N\!-\!n$ (notation below), and the relay $01$ is an edge $O\!-\!n$.

\begin{proposition}[computation]\label{prop:hubtest}
The following holds for the game with path insertion, receipt or flower, and the orbit rule
(\texttt{scripts/hs\_test.py}).
\begin{enumerate}[label=(\alph*)]
\item \emph{Order $8$.} The $11117$ connected graphs of order $8$ have $9268$ automorphism orbits
of roots satisfying the hypotheses of Conjecture~\ref{conj:hub}. With $F=\mathrm{star}(y)\cup X$,
the least $|X|$ is $0$ for $3100$ of them, $1$ for $5375$, $2$ for $787$ and $3$ for $6$, both
with the target of (R) and with the simple target. Always $\mathrm{lb}(y)\le1$, and
$|X|=\mathrm{lb}(y)$ for $8475$ root orbits, $\mathrm{lb}(y)+1$ for $783$ and $\mathrm{lb}(y)+2$
for $10$. So Conjecture~\ref{conj:hub} holds at order $8$ with $c=2$. If the star may be centred at any
vertex, the least number of further edges is $0$, $1$ or $2$ for $7000$, $2256$ and $12$ root
orbits ($6997$, $2259$ and $12$ with the simple target).
\item \emph{The proved theorems alone.} Theorem~\ref{thm:relayhub} applies at $2556$ of these
$9268$ root orbits, Corollary~\ref{cor:twohub} at $379$, and one of them at $2757$. At each of
them the game, with path insertion alone, wins with the theorem's own held set
(\texttt{scripts/relay\_cover.py}).
\item \emph{Order $10$.} We take the ten critical graphs of Proposition~\ref{prop:rfgame} and $40$
more sampled in the same way (\texttt{scripts/sample10.py}, seed $23$): critical connected graphs
of order $10$ with a vertex of degree $6$. Thirty of the $50$ have a root satisfying the
hypotheses, with $108$ root orbits in all, and $\mathrm{lb}(y)=0,1,2$ for $2$, $39$ and $67$ of
them. With $F=\mathrm{star}(y)\cup X$ and the target of (R), the least $|X|$ is $0$, $1$, $2$, $3$,
$4$ for $1$, $35$, $68$, $3$, $1$ root orbits. Then $|X|-\mathrm{lb}(y)$ is $0$ for $102$ root
orbits, $1$ for $3$, $2$ for $1$ and $3$ for $2$. The two with $|X|=\mathrm{lb}(y)+3$ are the roots
$7$ ($d=8$, $\mathrm{lb}=0$, $|X|=3$) and $9$ ($d=6$, $\mathrm{lb}=1$, $|X|=4$) of
\texttt{Inv\char126|\char126zbw}. The simple target gives the same values, except at the root $2$
of \texttt{Inv\char126|\char126zbw}, which needs $|X|=4=\mathrm{lb}(y)+2$ instead of $3$. If the star
may be centred at any vertex, the least number of further edges is $0$, $1$, $2$, $3$ for $3$,
$93$, $10$, $2$ root orbits, with either target. Neither Theorem~\ref{thm:relayhub} nor
Corollary~\ref{cor:twohub} applies at any of the $108$ root orbits.
\end{enumerate}
\end{proposition}
\begin{sloppypar}
The files in \texttt{scripts/outputs/} are \texttt{hs8\_R.log}, \texttt{hs8\_simple.log},
\texttt{hs8\_R\_root.log}, \texttt{hs8\_simple\_root.log}, \texttt{hs8\_root\_parity.txt} and
\texttt{relay\_cover8.txt} for order $8$. For order $10$ they are \texttt{hs10\_R.log},
\texttt{hs10\_simple.log}, \texttt{hs10\_R\_root.log}, \texttt{hs10\_simple\_root.log},
\texttt{hs10\_root\_parity.txt}, \texttt{hs10\_root\_parity\_simple.txt}, \texttt{relay\_cover10.txt}
and \texttt{relay\_cover10\_reasons.txt}. The graphs are in \texttt{sample10\_50.g6}.
\end{sloppypar}

The shapes of $X$ at order $8$ are small parity repairs (\texttt{scripts/hs\_shapes.py}, output
\texttt{hs8\_root\_shapes.txt}). Write $N$ and $O$ for an
even vertex in or outside $N(y)$, and $n$, $o$ for an odd one. The solver reports the first
winning $X$ in lexicographic order, so these counts describe one winning shape per root, not all
of them. For $|X|=1$ the edge is a relay $O\!-\!n$ in $2468$ cases, joins two odd neighbours in
$1495$, joins two vertices $O$ in $1332$, and is of another type in $80$. For $|X|=2$ it is a path
of length two in $757$ cases and a matching in $30$. Of the paths, $611$ have an even neighbour
of $y$ as the middle vertex, which then carries three held edges, like the second hub of
\texttt{Inv\char126|\char126zbw}. The other $146$ have an odd vertex outside $N[y]$ as the middle
vertex. For $|X|=3$ there are four paths of length three and two triangles.

\subsection*{Negative results}
\begin{enumerate}[label=(\arabic*)]
\item The bound $|X|\le\mathrm{lb}(y)+2$ fails at order $10$, at the roots $7$ and $9$ of
\texttt{Inv\char126|\char126zbw} (Proposition~\ref{prop:hubtest}(c)). The value $c=3$ is not
excluded by our data. The sample is small, though, and this graph is the only one in it that
needs more than $\mathrm{lb}(y)+1$. It is a dense graph, and (R) at its root $9$ follows
from~\cite[Thm.~12.29]{SP} (see after Proposition~\ref{prop:rfgame}).
\item The proved theorems are far from the computed wins. They cover $2757$ of the $9268$ root
orbits of order $8$, and none of the $108$ of order $10$. At order $10$ the relay theorem fails for
lack of even partners outside $N[y]$ for the odd neighbours ($99$ root orbits), by~(iii) ($4$)
or by~(ii) ($5$). Corollary~\ref{cor:twohub} needs $d(y)\le4$, which only $3$ of the $108$ have.
\item The game does not win the order-$10$ roots by the receipts-only pattern of
Lemma~\ref{lem:genhub}. At a root of degree $8$ a pure-star seed has $b\le1$ on $N(y)$, so
Theorem~\ref{thm:relayhub} would need $q\ge3$ relays. Yet $54$ of the $57$ root orbits of degree
$8$ win with $|X|=\mathrm{lb}(y)\le2$. Their plays must use other moves, such as donations from
the hub or receipt or flower, which our lemmas do not cover.
\end{enumerate}

\section{Consequences}\label{sec:conseq}

\begin{enumerate}
\item The conclusion of connections~\S R needs correcting. Lov\'asz seeds with Fan transfers were
incomplete at order $8$ only because path insertion was missing, and path insertion is already in
the Lean insertion library (\texttt{Ach.addPath}). With it, the game is complete at order $8$.
\item Inductive seeds, the rooted test on a connected spanning $G-F$, have $\sum b=n$, so they have
no spare path. The spare path comes from the floor $\lfloor N/2\rfloor$ at odd order, i.e.\ from
Lov\'asz's theorem with an isolated vertex. An induction that keeps it needs floor bounds at odd
order. These hold only in special classes: Lov\'asz's (at most one even vertex), and the bounds
$\mathrm{BP}_t$ of the roadmap (a pendant leaf and $t$ even vertices). Lemma~\ref{lem:hub} turns any
such floor partition of $G-s$, with an end at each neighbour of $s$, into (R) at a degree-four
vertex $s$.
\item P\'osa moves at the profile level (rotation at equality, Fan's alternative) repair no
losing root at order $8$. The difficulty is the degree-six hub, where a pure star is one receipt short at Fan's
equality. This makes the ``two-end transfer at equality'' of~\cite{D6} concrete. The rigid
configuration of Lemmas~\ref{lem:alt} and~\ref{lem:through} must be refuted by an edge count
(Proposition~\ref{prop:A62}) or recombined as a flower (Proposition~\ref{prop:K8M}). The count of
Lemma~\ref{lem:rfcount} makes both available in the game. It closes the hard cases $A_{6,2}$,
$B'$ and $\mathrm{CP}(4),\mathrm{CP}(5)$. It is not sharp enough on \texttt{Inv\char126|\char126zbw},
where a second hub along the inserted path closes the case instead.
Otherwise the configuration is bypassed by holding more edges and inserting a longer path through
the hub.
\item Lemmas~\ref{lem:receipts}--\ref{lem:twohub} are the receipts-only part of the hub calculus
of~\cite{D6} (receipts, donations, paybacks, batteries), with the spare path of~(A) as the payment.
They explain the order-$8$ certificates of the simplest shapes, but not the order-$10$ wins
(Section~\ref{sec:general}, negative results). A proof of \texttt{BranchBase3} along these lines
needs the donation and payback steps as well, and a simple final profile.
\end{enumerate}

\section{Formalisation and scripts}\label{sec:scripts}
\sloppy

\subsection*{Lean} File \texttt{lean/GallaiLean/MissingMove.lean} of the project repository (Lean~4,
Mathlib v4.36.0). It has no \texttt{sorry}, and \texttt{\#print axioms} reports only
\texttt{propext}, \texttt{Classical.choice} and \texttt{Quot.sound}.
\begin{itemize}
\item \texttt{Ach.hub\_four}: Lemma~\ref{lem:hub}, for the bookkeeping predicate
\texttt{Ach E b} (the graph with edge set $E$ has a partition with end vector $b$). The profile
identity is stated without subtraction:
$b'(v)+[v=a]+[v=y]=b(v)+2[v=s]+[v=y']+[v=x]$.
\item \texttt{hub\_four\_ends}, \texttt{hub\_four\_rooted}: Theorem~\ref{thm:hub4}, with
Lov\'asz's theorem from the library (\texttt{lovasz\_ends}).
\item Helpers: \texttt{isolate} ($G$ with every edge at $s$ removed), \texttt{isolate\_degree},
\texttt{isolate\_edgeSet}, and \texttt{Ach.decomp\_of\_edgeSet}, \texttt{Ach.of\_decomp}.
\item \texttt{Ach.hub\_four\_tail}, \texttt{hub\_four\_odd\_ends}, \texttt{hub\_four\_odd\_rooted}:
Theorem~\ref{thm:hub4odd}. \texttt{gqz\_rooted\_three}: Corollary~\ref{cor:gqz}, with the
hypotheses checked by \texttt{decide}.
\item \texttt{Ach.hub\_receipts}, \texttt{Ach.hub\_general}, \texttt{Ach.relays},
\texttt{Ach.relay\_hub} and \texttt{Ach.two\_hub}: Lemmas~\ref{lem:receipts}, \ref{lem:genhub},
\ref{lem:relays}, the step of Theorem~\ref{thm:relayhub} after the seed, and
Lemma~\ref{lem:twohub}. Profiles are again stated without subtraction, e.g.\ for
\texttt{Ach.hub\_general}:
$b'(v)+[v\in R]=b(v)+(d-2)[v=s]+[v=y]+[v=y']$.
\item \texttt{fan\_alternative}: Lemma~\ref{lem:alt}. Its two halves are
\texttt{fan\_reroute\_exec} and \texttt{fan\_alt\_count}. The first executes a partial chain; it is
the reroute part of \texttt{fan\_exec} from \texttt{Fan.lean}, without the final append. The second
is the counting of \texttt{fan\_exists\_chain}, stopped before the contradiction.
\end{itemize}
Lemma~\ref{lem:rot} is the one-step case of \texttt{fan\_reroute\_exec} together with
\texttt{PathDecomp.append\_move}. Proposition~\ref{prop:A62}, Lemmas~\ref{lem:flowermove}
and~\ref{lem:rfcount} are not formalised. The emc-four session's file
\texttt{lean/GallaiLean/Flower.lean} formalises Lemma~\ref{lem:flower} (\texttt{flower\_decomp}),
Proposition~\ref{prop:K8M} (\texttt{cocktail\_passes\_four}) and the uniform reduction of
$\mathrm{CP}(m)$ to a transversal lemma (\texttt{cocktail\_passes\_of\_rt}). Its transversal
lemmas \texttt{rainbowTransversal\_three} and \texttt{rainbowTransversal\_four} are the cases
$j=3$ on $8$ vertices and $j=4$ on $10$ vertices, in the weaker form of a path \emph{through} $u$.
Lemma~\ref{lem:transversal} in the form used here, a path \emph{starting} at $u$, is not
formalised.
The Lov\'asz wrappers of Theorem~\ref{thm:relayhub} and Corollary~\ref{cor:twohub} are not formalised either.

\subsection*{Scripts} In \texttt{scripts/}:
\begin{itemize}
\item \texttt{fan\_game.py}: the game solver, a copy of \texttt{tools/fan\_game.py}. Its options
are \texttt{np} (path insertion), \texttt{rotR} (at most $R$ rotations), \texttt{altR} (at most $R$
uses of Fan's alternative), \texttt{ind} (inductive seeds), \texttt{hubs} (held sets of the
form ``full star plus $k$ edges''), \texttt{hubroot} (the same, with the star at the root),
\texttt{flower} (Lemma~\ref{lem:flowermove}), \texttt{rf} (Lemma~\ref{lem:rfcount}), \texttt{orbit} (the
orbit rule) and \texttt{simple} (the simple target). The split rule is always on; it
does not change the results of Section~\ref{sec:eight}.
\item \texttt{orbit\_run.py}, \texttt{single\_root.py}, \texttt{fam\_hub34.py},
\texttt{miss\_k3.py}: runs at one root per automorphism orbit, at a single root and held-set size,
for the order-$10$ families with a star plus three edges, and for \texttt{Inv\char126|\char126zbw}.
\item \texttt{crit.py}, \texttt{sample10.py}, \texttt{rf\_sample.sh}: criticality, the random
order-$10$ sample, and the sample run of Proposition~\ref{prop:rfgame}.
\item \texttt{certificate.py}: prints a winning strategy tree, with path insertion, for one graph and
root.
\item \texttt{stepper.sh}: for each graph, raises $|F|$ from $5$ to $8$ until every root wins.
\item \texttt{families.py}: the order-$10$ test graphs $A_{4,4}$, $A_{6,2}$, $B$, $B'$.
\item \texttt{rootcheck.py}, with \texttt{pathnum2.py}: an independent CP-SAT check of (R).
\item \texttt{hs\_test.py}: the test of Conjecture~\ref{conj:hub}, per root orbit (options
\texttt{hubroot}: star at the root; \texttt{simple}: simple target). \texttt{hs\_parity.py} compares
the least $|X|$ with $\mathrm{lb}(y)$, \texttt{hs\_shapes.py} classifies $X$, and
\texttt{relay\_cover.py} checks the hypotheses of Theorem~\ref{thm:relayhub} and
Corollary~\ref{cor:twohub} and replays their held sets. \texttt{sample10.py} samples critical
graphs of order $10$.
\item \texttt{run\_all.sh}: reproduces every table. The outputs are in \texttt{scripts/outputs/}.
The graph lists (\texttt{.g6}) and certificates (\texttt{certificate*.txt}) are committed. The
\texttt{.log} files that \texttt{run\_all.sh} writes, such as \texttt{crit47.log},
\texttt{famhub\_B2\_3.log} and \texttt{rf\_sample.log}, are excluded by the repository's
\texttt{.gitignore} and must be regenerated.
\end{itemize}

\begin{thebibliography}{99}
\bibitem{L68} L.~Lov\'asz, On covering of graphs, in: \emph{Theory of Graphs} (Proc.\ Colloq.\
Tihany 1966), Academic Press, 1968, 231--236.
\bibitem{Fan} G.~Fan, Path decompositions and Gallai's conjecture, \emph{J.\ Combin.\ Theory Ser.\
B} 93 (2005) 117--125.
\bibitem{ES} J.~Williams, Endpoint sets of path decompositions at the Gallai budget, working draft,
\texttt{manuscripts/endpoint-structure}.
\bibitem{D6} The $d=6$ frontier: transfers, relays and the terminal step, working draft,
\texttt{manuscripts/d6-frontier}.
\bibitem{EXP} Project experiment log, \texttt{research/experiments.md}, entry E13.
\bibitem{SP} J.~Williams, Separators and endpoint exchange in the rooted Gallai problem, working
draft, 25 November 2026, \texttt{gallai/02-november-25/separating-pairs.pdf}.
\bibitem{MR} J.~Williams, Cut vertices and cycle intersections in Gallai's path decomposition
problem, working draft, 24 November 2026,
\texttt{gallai/01-original-manuscripts/marked-reduction.pdf}.
\bibitem{CC} J.~Williams, Nearly spanning cycles and long intersections in Gallai's path
decomposition problem, working draft, 26 November 2026,
\texttt{gallai/03-november-26-papers/cycle-and-contact.pdf}.
\bibitem{JM} R.~Javadi, M.~Miralaei, A conjecture on rainbow Hamiltonian cycle decomposition,
\emph{Discrete Math.} 349 (2026) 114676; arXiv:2403.17290.
\bibitem{JMR} Repairing the extension step of Javadi--Miralaei, working draft,
\texttt{manuscripts/jm-repair}.
\end{thebibliography}

\end{document}
